\paragraph{Fourier conductor 壳层、真商模下降与 full-modulus 原始性}
在折叠多重度剖面的 Parseval 恒等式、群环乘积分解以及 Fibonacci 共振双频的常数级下界已经建立之后，还可把“是否能够下降到真商模”完全转写为 Fourier 支撑的 conductor 截断问题。

\begin{theorem}[折叠多重度剖面下降到真商模的 Fourier--conductor 充要条件]\label{thm:xi-fold-multiplicity-quotient-descent-fourier-conductor}
记
\[
M:=F_{m+2},
\qquad
d_m:\ZZ/(M\ZZ)\to\NN,
\qquad
\widehat d_m(k):=\sum_{r\in\ZZ/(M\ZZ)}d_m(r)e^{-2\pi i kr/M}.
\]
固定约数 \(D\mid M\)。则下列三条命题等价：
\begin{enumerate}
  \item 存在函数
  \[
  f_D:\ZZ/(D\ZZ)\to\NN
  \]
  使得
  \[
  d_m(r)=f_D(r\bmod D)
  \qquad
  (\forall r\in\ZZ/(M\ZZ));
  \]
  \item 对每个频率 \(k\in\ZZ/(M\ZZ)\)，只要
  \[
  k\not\equiv 0\pmod{M/D},
  \]
  就有
  \[
  \widehat d_m(k)=0;
  \]
  \item 若定义导出模
  \[
  \operatorname{cond}(k):=\frac{M}{\gcd(k,M)}
  \]
  及各壳层能量
  \[
  E_m(q):=\sum_{\operatorname{cond}(k)=q}|\widehat d_m(k)|^2
  \qquad (q\mid M),
  \]
  则
  \[
  E_m(q)=0
  \qquad
  \text{对一切 }q\nmid D.
  \]
\end{enumerate}
此外，不加任何额外假设时总有精确分解
\[
\mathsf{Col}_m
=
\frac1M+\frac{1}{2^{2m}M}\sum_{\substack{q\mid M\\ q>1}}E_m(q),
\]
而在上述等价条件成立时，上式退化为
\[
\mathsf{Col}_m
=
\frac1M+\frac{1}{2^{2m}M}\sum_{\substack{q\mid D\\ q>1}}E_m(q).
\]
\end{theorem}
\begin{proof}
先证 \((1)\Rightarrow(2)\)。若
\[
d_m(r)=f_D(r\bmod D),
\]
则把 \(r\) 唯一写成
\[
r=a+Db,
\qquad
0\le a<D,
\qquad
0\le b<\frac{M}{D}.
\]
于是
\[
\begin{aligned}
\widehat d_m(k)
&=
\sum_{a=0}^{D-1}\sum_{b=0}^{M/D-1}
f_D(a)\,e^{-2\pi i k(a+Db)/M}\\
&=
\left(\sum_{a=0}^{D-1}f_D(a)e^{-2\pi i ka/M}\right)
\left(\sum_{b=0}^{M/D-1}e^{-2\pi i kb/(M/D)}\right).
\end{aligned}
\]
第二个括号是长度 \(M/D\) 的几何和；当且仅当
\[
\frac{M}{D}\mid k
\]
时非零，否则为零。故 \((2)\) 成立。

再证 \((2)\Rightarrow(1)\)。Fourier 反演给出
\[
d_m(r)=\frac1M\sum_{k\in\ZZ/(M\ZZ)}\widehat d_m(k)e^{2\pi i kr/M}.
\]
在 \((2)\) 的假设下，只保留 \(k=(M/D)\ell\) 的项，可得
\[
\begin{aligned}
d_m(r)
&=
\frac1M\sum_{\ell=0}^{D-1}\widehat d_m\!\left(\frac{M}{D}\ell\right)
e^{2\pi i \ell r/D}.
\end{aligned}
\]
右端仅依赖于 \(r\bmod D\)，故 \(d_m\) 下降到 \(\ZZ/(D\ZZ)\)。

现证 \((2)\Leftrightarrow(3)\)。若
\[
k=\frac{M}{D}\ell,
\]
则
\[
\gcd(k,M)=\frac{M}{D}\gcd(\ell,D),
\qquad
\operatorname{cond}(k)=\frac{D}{\gcd(\ell,D)}\mid D.
\]
反之，若 \(\operatorname{cond}(k)\mid D\)，记
\[
g:=\gcd(k,M),
\qquad
\operatorname{cond}(k)=\frac{M}{g}.
\]
由 \(\frac{M}{g}\mid D\) 且 \(D\mid M\)，可写
\[
D=\frac{M}{g}\,t
\]
对某个整数 \(t\)。于是
\[
\frac{M}{D}=\frac{g}{t},
\]
从而 \(\frac{M}{D}\mid g\)。由于 \(g\mid k\)，遂有
\[
\frac{M}{D}\mid k.
\]
故
\[
\frac{M}{D}\mid k
\iff
\operatorname{cond}(k)\mid D,
\]
于是 \((2)\) 与 \((3)\) 等价。

最后证明碰撞分解。由定理 \ref{thm:fold-collision-spectrum}，
\[
\mathsf{Col}_m
=
\frac{1}{2^{2m}M}\sum_{k\in\ZZ/(M\ZZ)}|\widehat d_m(k)|^2.
\]
其中 \(k=0\) 一项恰给出 \(1/M\)，而非零频率按 \(\operatorname{cond}(k)\) 分组，便得到
\[
\mathsf{Col}_m
=
\frac1M+\frac{1}{2^{2m}M}\sum_{\substack{q\mid M\\ q>1}}E_m(q).
\]
若 \((1)\)--\((3)\) 成立，则所有 \(q\nmid D\) 的壳层均为零，故立即退化为
\[
\mathsf{Col}_m
=
\frac1M+\frac{1}{2^{2m}M}\sum_{\substack{q\mid D\\ q>1}}E_m(q).
\]
证毕。
\end{proof}
\leanverified{paper\_xi\_fold\_multiplicity\_quotient\_descent\_fourier\_conductor}

\begin{corollary}[大尺度下折叠多重度剖面的 full-modulus 原始性]\label{cor:xi-fold-multiplicity-fullmodulus-primitive}
存在整数 \(m_0\)，使得对所有 \(m\ge m_0\)，折叠多重度剖面
\[
d_m:\ZZ/(F_{m+2}\ZZ)\to\NN
\]
都不可能下降到任何真商模
\[
\ZZ/(D\ZZ),
\qquad
D\mid F_{m+2},
\qquad
D<F_{m+2}.
\]
等价地，对充分大的 \(m\)，其最小可见模数正是 \(F_{m+2}\) 本身。
\end{corollary}
\begin{proof}
由定理 \ref{thm:fold-critical-resonance-constant}，
\[
a_m(F_m):=\frac{|\widehat d_m(F_m)|}{2^m}\longrightarrow C_\varphi>0.
\]
故存在 \(m_0\)，使得当 \(m\ge m_0\) 时，
\[
\widehat d_m(F_m)\neq 0.
\]
另一方面，由 Fibonacci 最大公因子恒等式，
\[
\gcd(F_m,F_{m+2})=F_{\gcd(m,m+2)}=1,
\]
从而
\[
\operatorname{cond}(F_m)=\frac{F_{m+2}}{\gcd(F_m,F_{m+2})}=F_{m+2}.
\]
若 \(d_m\) 能下降到某个真商模 \(D<F_{m+2}\)，则由定理
\ref{thm:xi-fold-multiplicity-quotient-descent-fourier-conductor} 的条件 \((3)\)，所有
\[
q\nmid D
\]
的 conductor 壳层都必须消失。特别地，full-conductor 壳层
\[
q=F_{m+2}
\]
必须为零，这与 \(\widehat d_m(F_m)\neq 0\) 矛盾。故不存在这样的真商模 \(D\)。证毕。
\end{proof}

\begin{theorem}[full-conductor 壳层的碰撞缺口下界由 Fibonacci 共振双频承担]\label{thm:xi-full-conductor-shell-collision-gap-fibonacci-pair}
记
\[
M:=F_{m+2},
\qquad
E_m^{\mathrm{full}}:=E_m(M)=\sum_{\gcd(k,M)=1}|\widehat d_m(k)|^2,
\qquad
a_m(k):=\frac{|\widehat d_m(k)|}{2^m}.
\]
则有
\[
E_m^{\mathrm{full}}
\ge
|\widehat d_m(F_m)|^2+|\widehat d_m(F_{m+1})|^2
=
2^{2m}\bigl(a_m(F_m)^2+a_m(F_{m+1})^2\bigr).
\]
从而
\[
\mathsf{Col}_m-\frac1M
\ge
\frac{E_m^{\mathrm{full}}}{2^{2m}M}
\ge
\frac{a_m(F_m)^2+a_m(F_{m+1})^2}{M}
=
\frac{2C_\varphi^2+o(1)}{M}.
\]
特别地，碰撞缺口的主尺度下界已经完全由 full-conductor 壳层内的共轭 Fibonacci 双频承担。
\end{theorem}
\begin{proof}
由定理 \ref{thm:xi-fold-multiplicity-quotient-descent-fourier-conductor} 的壳层分解，
\[
\mathsf{Col}_m-\frac1M
=
\frac{1}{2^{2m}M}\sum_{\substack{q\mid M\\ q>1}}E_m(q)
\ge
\frac{E_m^{\mathrm{full}}}{2^{2m}M}.
\]

又由 Fibonacci 最大公因子恒等式，
\[
\gcd(F_m,F_{m+2})=\gcd(F_{m+1},F_{m+2})=1,
\]
故频率 \(F_m\) 与 \(F_{m+1}\) 都落在 full-conductor 壳层内。于是
\[
E_m^{\mathrm{full}}
\ge
|\widehat d_m(F_m)|^2+|\widehat d_m(F_{m+1})|^2.
\]
再用
\[
|\widehat d_m(F_m)|=2^m a_m(F_m),
\qquad
|\widehat d_m(F_{m+1})|=2^m a_m(F_{m+1}),
\]
便得到第二个不等式。

最后，由定理 \ref{thm:fold-critical-resonance-constant}，
\[
a_m(F_m)\to C_\varphi,
\qquad
a_m(F_{m+1})\to C_\varphi,
\]
故
\[
\frac{a_m(F_m)^2+a_m(F_{m+1})^2}{M}
=
\frac{2C_\varphi^2+o(1)}{M}.
\]
证毕。
\end{proof}

\begin{theorem}[Fourier 幅度的中点门级联实现与临界共振半径]\label{thm:xi-time-part9p1-midpoint-gate-cascade-critical-radius}
沿用定理 \ref{thm:fold-spectrum-factorization} 的记号。对每个 \(m\ge 1\) 与 \(k\in\ZZ/(F_{m+2}\ZZ)\)，定义角序列
\[
\vartheta_j^{(m)}(k):=
2\pi k\frac{F_{j+1}}{F_{m+2}}
\qquad (1\le j\le m),
\]
以及复平面上的中点门
\[
\mathcal M_{\vartheta}(z):=\frac{z+e^{-\mathrm{i}\vartheta}z}{2}.
\]
再置
\[
z_0^{(m,k)}:=1,
\qquad
z_j^{(m,k)}:=\mathcal M_{\vartheta_j^{(m)}(k)}\bigl(z_{j-1}^{(m,k)}\bigr)
\qquad (1\le j\le m).
\]
则成立：
\begin{enumerate}
  \item 级联终值满足
  \[
  \boxed{\;
  z_m^{(m,k)}
  =
  2^{-m}\widehat d_m(k),\;}
  \]
  从而
  \[
  \boxed{\;
  |z_m^{(m,k)}|=a_m(k).\;}
  \]
  因此 \(a_m(k)\) 精确等同于角序列 \(\{\vartheta_j^{(m)}(k)\}_{j=1}^{m}\) 所驱动的中点门级联收缩半径；
  \item 有
  \[
  \widehat d_m(k)=0
  \iff
  z_m^{(m,k)}=0
  \iff
  \exists\, j\in\{1,\dots,m\}\ \text{使}\ 
  \vartheta_j^{(m)}(k)\equiv \pi\pmod{2\pi}.
  \]
  换言之，Fourier 零点恰对应于某一级中点门把一点与其对径旋转副本取平均而坍缩到原点；
  \item 对 Fibonacci 补频对 \(k=F_m\) 与 \(k=F_{m+1}\)，有
  \[
  |z_m^{(m,F_m)}|\longrightarrow C_\varphi,
  \qquad
  |z_m^{(m,F_{m+1})}|\longrightarrow C_\varphi.
  \]
  因而临界共振常数 \(C_\varphi\) 可解释为一条中点门级联在 Fibonacci 角驱动下保持的严格正极限半径。
\end{enumerate}
\end{theorem}
\begin{proof}
由定义，
\[
\mathcal M_{\vartheta}(z)=\frac{1+e^{-\mathrm{i}\vartheta}}{2}\,z.
\]
故递推展开给出
\[
z_m^{(m,k)}
=
\prod_{j=1}^{m}\frac{1+e^{-\mathrm{i}\vartheta_j^{(m)}(k)}}{2}
=
2^{-m}\prod_{j=1}^{m}\left(1+e^{-2\pi \mathrm{i}kF_{j+1}/F_{m+2}}\right).
\]
再由定理 \ref{thm:fold-spectrum-factorization}，
\[
\widehat d_m(k)=\prod_{j=1}^{m}\left(1+e^{-2\pi \mathrm{i}kF_{j+1}/F_{m+2}}\right),
\]
便得到
\[
z_m^{(m,k)}=2^{-m}\widehat d_m(k).
\]
取模即得
\[
|z_m^{(m,k)}|=\frac{|\widehat d_m(k)|}{2^m}=a_m(k).
\]

第二条由乘积分解立即推出：\(\widehat d_m(k)=0\) 当且仅当某一因子
\[
1+e^{-\mathrm{i}\vartheta_j^{(m)}(k)}
\]
为零，而这又等价于
\[
\vartheta_j^{(m)}(k)\equiv \pi\pmod{2\pi}.
\]
此时
\[
\mathcal M_{\vartheta_j^{(m)}(k)}(z)=0
\qquad (\forall z\in\CC),
\]
故级联在该步后永久坍缩到原点。

第三条正是定理 \ref{thm:fold-critical-resonance-constant} 的结论，因为
\[
|z_m^{(m,F_m)}|=a_m(F_m),
\qquad
|z_m^{(m,F_{m+1})}|=a_m(F_{m+1}).
\]
证毕。
\end{proof}

\endinput
